% META: {"id": "TNSI-LANG-M1", "chapitre": "TNSI-LANGAGES-ET-PROGRAMMATION", "type_objet": "methode", "capacites": ["T-LANG-02A", "T-LANG-02B"], "capacites_codes": ["C4", "C5"], "methodes": ["M1"], "mode_creation": "ex_nihilo", "sources_inspiration": [], "status": "needs_review"}
\begin{fichemethode}{M1}{Écrire une fonction récursive et dérouler son arbre d'appels}
\textbf{Quand l'utiliser ?} Dans les deux sens : quand on te demande d'écrire une fonction
récursive sur un problème neuf, et quand on te donne une fonction récursive et qu'on demande
ce qu'elle renvoie, combien d'appels elle fait, ou si elle termine. Écrire et dérouler sont la
même démarche prise par les deux bouts.

\textbf{Pas à pas — écrire :}
\begin{enumerate}
  \item \textbf{Écris le cas de base en premier}, avant toute autre ligne. Quelle est la plus
    petite entrée pour laquelle tu connais la réponse sans réfléchir ? La liste vide, la
    chaîne vide, $n = 0$.
  \item \textbf{Écris comment le cas général se RAMÈNE à plus petit.} Une seule règle compte :
    l'argument de l'appel récursif doit être \textbf{strictement} plus proche du cas de base.
    Note cette quantité qui décroît — c'est ta preuve de terminaison.
  \item \textbf{Suppose l'appel récursif correct} et écris la ligne qui combine son résultat
    avec le travail du niveau courant. C'est la ligne la plus courte et la plus difficile.
  \item \textbf{Relis la terminaison} : est-ce que tout chemin d'exécution atteint le cas de
    base ? Une branche qui rappelle la fonction avec le même argument suffit à tout casser.
\end{enumerate}

\textbf{Pas à pas — dérouler :}
\begin{enumerate}
  \item \textbf{Écris l'appel initial}, puis, sous lui et décalés vers la droite, les appels
    qu'il déclenche, dans l'ordre du code.
  \item \textbf{Descends jusqu'aux cas de base} : ce sont les feuilles de l'arbre, les seules
    qui renvoient sans rappeler.
  \item \textbf{Remonte en écrivant les valeurs de retour}, feuille par feuille, en appliquant
    la ligne de combinaison.
  \item \textbf{Compte.} Le nombre d'appels est le nombre de nœuds de l'arbre. Repère si un
    même argument apparaît dans deux branches : c'est le signe qu'une mémoïsation ferait
    gagner (voir la fiche du chapitre Algorithmique).
\end{enumerate}

\textbf{Exemple :} \lstinline{somme_chiffres(n)}, somme des chiffres d'un entier positif.

\emph{Cas de base} : si $n < 10$, la somme est $n$ lui-même.

\emph{Réduction} : $n$ vaut \lstinline{n // 10} avec un chiffre en moins ; la quantité qui
décroît strictement est le nombre de chiffres.

\emph{Combinaison} : \lstinline{n % 10 + somme_chiffres(n // 10)}.

\emph{Déroulé de} \lstinline{somme_chiffres(507)} : $7 + \text{s}(50)$, puis
$0 + \text{s}(5)$, puis $5$ — cas de base. On remonte : $5$, puis $0+5=5$, puis $7+5=12$.
Trois appels, un seul par niveau : l'arbre est une chaîne, aucune mémoïsation à espérer.

% BEGIN-VERIFY
% def somme_chiffres(n):
%     if n < 10:
%         return n
%     return n % 10 + somme_chiffres(n // 10)
%
% # Correction contre un oracle independant.
% for n in range(0, 3000):
%     assert somme_chiffres(n) == sum(int(c) for c in str(n))
% assert somme_chiffres(507) == 12
%
% # Etape 2 : la quantite qui decroit est le nombre de chiffres, strictement.
% for n in range(10, 5000):
%     assert len(str(n // 10)) < len(str(n))
%
% # Le deroule annonce : trois appels pour 507, un seul par niveau.
% trace = []
% def somme_tracee(n):
%     trace.append(n)
%     if n < 10:
%         return n
%     return n % 10 + somme_tracee(n // 10)
% assert somme_tracee(507) == 12
% assert trace == [507, 50, 5]
% assert len(trace) == len(set(trace)) == 3   # aucun argument repete
%
% # Contraste : un arbre d'appels qui SE RECOUPE, ou la memoisation servirait.
% appels = {}
% def fib(n):
%     appels[n] = appels.get(n, 0) + 1
%     if n <= 1:
%         return n
%     return fib(n - 1) + fib(n - 2)
% assert fib(5) == 5
% assert appels[3] == 2 and appels[2] == 3    # memes arguments, plusieurs branches
% assert sum(appels.values()) == 15           # 15 noeuds dans l'arbre d'appels
%
% # Etape 4 : une recursion qui ne reduit pas ne termine pas. On le montre avec
% # une profondeur bornee plutot qu'en laissant Python lever RecursionError.
% def sans_reduction(n, profondeur=0):
%     if profondeur > 60:
%         return "NON_TERMINE"
%     if n < 10:
%         return n
%     return sans_reduction(n, profondeur + 1)   # n ne change jamais
% assert sans_reduction(507) == "NON_TERMINE"
%
% # Un grand nombre d'appels n'est PAS une non-terminaison : fib termine.
% assert fib(18) == 2584
% END-VERIFY

\textbf{Pièges :}
\begin{itemize}
  \item \textbf{Écrire le cas général avant le cas de base.} C'est l'ordre qui fait oublier le
    cas de base, et sans lui rien ne s'arrête.
  \item \textbf{Un appel récursif qui ne réduit rien.} \lstinline{f(n)} qui rappelle
    \lstinline{f(n)} ne termine jamais, même si le cas de base existe. Nomme toujours la
    quantité qui décroît.
  \item \textbf{Confondre « beaucoup d'appels » et « ne termine pas ».}
    \lstinline{fib(30)} fait plus de deux millions d'appels et termine parfaitement. Le
    problème est le coût, pas la terminaison, et la réponse n'est pas la même.
  \item \textbf{Dérouler dans le désordre.} L'arbre d'appels se lit dans l'ordre du code :
    l'appel écrit à gauche se déclenche en premier, et sa valeur revient en premier.
\end{itemize}

\textbf{Vérifier :} teste ta fonction sur le cas de base, sur le cas juste au-dessus, et sur
un cas moyen comparé à une méthode directe. Pour un déroulé, compte les feuilles : elles
doivent toutes être des cas de base, et aucun nœud interne ne doit avoir la même valeur de
retour qu'un de ses enfants sans qu'une raison l'explique.

\textbf{S'entraîner :} exercices C4 et C5 du chapitre.
\end{fichemethode}
